\clearpage
\begingroup
\fancyhead[L]{\small BEMERKUNGEN}
\renewcommand{\headrulewidth}{0pt}
\par\noindent\hspace*{\parindent}Zu [VII.]: Die Briefstellen sind nach den Originalen abgedruckt, während
Notiz [2.] gänzlich vom Bearbeiter herrührt. Die ungeheuer umfangreichen
Rechnungen zu diesen Tafeln (vgl. S. 429), sowie die fertige Tafel selbst, bilden
ein Convolut von wohlgeordneten Blättern; viele von ihnen sind nicht von
\textsc{Gauss}' Hand, sondern rühren von \textsc{Encke} und \textsc{Westphal} her.

\par\noindent\hspace*{\parindent}Zu [VIII.]: Über die Störungen der Pallas durch Saturn befindet sich im
Nachlass nichts von \textsc{Gauss}' Hand, sondern nur die Originalbriefe von
\textsc{Nicolai} an \textsc{Gauss}, sowie einige Blätter von \textsc{Nicolai}s Hand;
aus den letztern stammen die Gleichungen 1) und 2) der Notiz [1.], während der
übrige Inhalt dieser Notiz ergänzt wurde. Die in Frage kommenden Briefstellen
sind in Notiz [2.]--[6.] abgedruckt; die Resultate für die Störungen stehen auf
einzelnen Blättern, die ursprünglich den \textsc{Nicolai}schen Briefen beigelegen
haben, jetzt aber im Nachlass den Kapseln mit den Pallasstörungen eingeordnet sind.
Übrigens sind diese Resultate in etwas anderer Ordnung abgedruckt, als sie auf
den Originalblättern stehen; dort sind sie in Manipel geordnet, während für den
Abdruck ihre ursprüngliche natürliche Reihefolge gewählt ist.

\par\noindent\hspace*{\parindent}Die ganze Rechnung dieser Störungen ist in \textsc{Nicolai}s Nachlass
auf der Grossherzoglichen Sternwarte zu Heidelberg erhalten und Herr
\textsc{Valentiner} hat die Freundlichkeit gehabt, sie dem Unterzeichneten zur
Einsicht zu überlassen. Sie trägt die Unterschrift: »Geendigt den
20\textsuperscript{ten} April 1814.« Vielleicht dürfen wir ihre etwas detaillirtere
Veröffentlichung, die in \textsc{Gauss}' Werken wohl nicht angebracht gewesen wäre,
noch von anderer Seite erhoffen. Übrigens hat \textsc{Nicolai} zwei Rechnungen der
Saturnstörungen ausgeführt, von denen die erste sich nur auf die Hauptgleichungen
bezog; ihre Resultate sind ebenfalls im Nachlass vorhanden. Die entsprechenden
Briefe von \textsc{Gauss} an \textsc{Nicolai} scheinen verloren zu sein; wenigstens
theilte mir Herr \textsc{Valentiner} mit, dass in \textsc{Nicolai}s Nachlass sich nur
eine Reihe von Briefen von \textsc{Gauss} aus späterer Zeit (von 1819 ab) vorfindet,
die Herr \textsc{Valentiner} übrigens herausgegeben hat (Briefe von C. F.
\textsc{Gauss} an B. \textsc{Nicolai}, Karlsruhe 1877). Die Briefe von \textsc{Nicolai}
an \textsc{Gauss} scheinen vollzählig erhalten zu sein, sind aber bisher nicht
veröffentlicht worden.

\par\noindent\hspace*{\parindent}Zu [IX.]: Die Formeln und Zahlen der Notiz [1.] finden sich theils im
Handbuch Bc, theils auf Blättern, theils in einem »Astronomische Rechnungen;
Angefangen im März 1814« betitelten Heft, das sonst andere Gegenstände enthält
und darum im Nachlass unter »Beobachtungen und Rechnungen« eingereiht ist.
Die Zahlen der Notizen [2.]--[3.] sind den von \textsc{Gauss} auf einzelnen Blättern
geführten Rechnungen entnommen, welche unvollendet geblieben sind. Die Notiz [4.]
stammt vom Bearbeiter; [5.]--[6.] aus dem erwähnten Heft »Astronomische Rechnungen«,
einzelne Relationen der Notiz [5.] auch von einem einzelnen Blatt; [7.] endlich
aus dem Handbuch Be, wo auf den abgedruckten Text noch ein Rechnungsbeispiel
(für \(M=165^0\)) folgt.

\par\noindent\hspace*{\parindent}Die Notizen [4.]--[7.] betreffen, wie man sieht, einen sehr interessanten
Kunstgriff für die Entwickelung der Störungsfunction, den \textsc{Gauss} aber
schliesslich doch nicht weiter angewandt hat, weil eine wesentliche Ersparniss
im numerischen Rechnen damit nicht erreicht zu werden scheint. Nach dieser
Methode hat \textsc{Gauss} ausser den abgedruckten Werthen für \(M=30^0\) nur noch
die für \(M=15^0\) und \(45^0\) berechnet.

\medskip
\par\noindent\hspace*{\parindent}Man sieht, dass \textsc{Gauss} das gewaltige Problem der Berechnung der
Pallasstörungen innerhalb der Genauigkeitsgrenze der Beobachtungen, an das auch
heute noch der Astronom sich nicht gern heranwagen möchte, fast ganz zu Ende
geführt hat; es fehlen nur einige wenige Störungsgleichungen für den Mars.
Um so mehr zu bedauern war es, dass bis auf den heutigen Tag diese Arbeit
unbekannt geblieben ist, und wie aus den letzten der abgedruckten Briefstellen
hervorgeht, hat \textsc{Gauss} selbst dies auch lebhaft gefühlt.

\par\noindent\hspace*{\parindent}Sehr auffallend ist, dass \textsc{Gauss} in seinem wissenschaftlichen
Tagebuch (oder Notizenjournal, wie er es selbst nennt), das von Herrn
F. \textsc{Klein} in der »Festschrift zur Feier des 150jährigen Bestehens der K. Ge%
\SourceLeaf{616}sellschaft der Wissenschaften zu Göttingen«, 1901 herausgegeben
worden ist, dieser seiner Untersuchung nicht ein einziges Mal Erwähnung thut.
Über die Jahre 1802--1804 steht im Tagebuch die Bemerkung: »\textit{Annis
insequentibus 1802, 1803, 1804 occupationes astronomicae maximam otii partem
abstulerunt, calculi imprimis circa planetarum novorum theoriam instituti.
Unde evenit, quod hisce annis catalogus hicce neglectus est} \mbox{.\ .\ .\ .\ .\ .\ .}«
In diese Jahre fallen \textsc{Gauss}' Arbeiten über das Zweikörperproblem
(Vorbereitung der Theoria motus) und über die Störungen der Ceres. Merkwürdig
ist es aber doch, dass \textsc{Gauss}' Beschäftigung mit astronomischen Rechnungen
stets eine Unterbrechung in der Führung des Tagebuchs hervorrief. Denn auch
im Jahre 1810, als \textsc{Gauss} begann, sich mit der Pallas zu beschäftigen,
beginnt eine grosse Lücke. Am 29. Februar 1812 schreibt er dann: »\textit{Catalogum
praecedentem per fata iniqua iterum interruptum initio anni 1812 resumimus}
\mbox{.\ .\ .\ .\ .}« Er erwähnt dann die Wiederauffindung des »zweiten« Beweises
des Fundamentalsatzes der Algebra, den er im November 1811 gefunden, dann
aber zum Theil wieder aus dem Gedächtniss verloren hatte, da er nichts
aufgezeichnet hatte. Zwischen dieser und der nächsten Notiz vom 26. September
1812, welche sich auf die Anziehung eines Ellipsoides bezieht, liegt die
Entdeckung der Rationalität der mittlern Bewegungen von Pallas und Jupiter;
auch über sie schweigt das Tagebuch vollständig.

\par\noindent\hspace*{\parindent}Dass die \textsc{Gauss}schen Untersuchungen über die Störungen der Pallas
noch heut ein werthvolles Material bilden, liegt auf der Hand; die Reduction
der Störungsausdrücke auf die heutige Epoche möchte man wohl als höchst
wünschenswerth bezeichnen.
\par\nobreak\noindent\hfill\textsc{Brendel}.\par
\nobreak\vspace{3pt}\noindent\makebox[\linewidth]{\rule{.24\linewidth}{.4pt}}\par
\clearpage
\endgroup
